\documentclass[11pt]{article}
\usepackage{fontspec}
\usepackage{amsmath,amssymb,amsthm,mathtools}
\usepackage[margin=1.1in]{geometry}
\usepackage{microtype}
\usepackage{parskip}
\usepackage{lmodern}
\usepackage[colorlinks=true,linkcolor=blue!50!black,urlcolor=blue!50!black]{hyperref}
\providecommand{\tightlist}{}
\title{A structural two-stage test for finite-state noise models, with results on public surface-code data}
\author{Felix Robles Elvira (ORCID: 0009-0009-2017-4394; independent researcher)}
\date{\small preprint, not peer reviewed, version 2026-06-11.}
\begin{document}
\maketitle

\section*{Abstract}
We present a structural (non-parametric) test of quantum-error- correction syndrome streams against entire classes of noise models at once, together with a cautionary correction to a natural first version of the test that we believe is of independent interest. \textbf{Stage one} (the sign stage): if a stationary binary stream is generated by a finite hidden-Markov model \textbf{in Moore form} (state-conditional emissions) whose underlying chain is \emph{reversible} (detailed balance), then the run statistics $m_k = p(1^{k+1})$ form a Hamburger moment sequence, so the offset Hankel matrices $[m_{i+j}]$ are positive semidefinite at every order — a theorem we prove by detailed-balance symmetrization.  Two scope corrections discipline the statement, and we present both because each is tempting to get wrong.  First, the $n = 0$ entry $p(1^0) = 1$ must be excluded: including it (the naive Hankel) falsely rejects legitimate periodic models, as a deterministic two-state alternator shows.  Second — found by hostile review of this paper's own earlier draft — the Moore qualifier is \emph{essential}: a Mealy (edge-emission) model on a reversible chain is \textbf{not} covered, because the induced Moore model lives on the edge chain, and the edge chain of a reversible chain is not reversible; indeed our own exact mechanism model (Lemma 2.3) \emph{is} a Mealy model on a reversible chain — i.i.d. measurement errors with symmetric edge emissions — and it signs. A statistically signed offset Hankel therefore certifies that \emph{no Moore-form hidden-Markov model with reversible chain, of any dimension,} reproduces the stream — a representation-class statement, narrower than the earlier draft's claim and stated exactly.  \textbf{Stage two} (the pole stage): Prony analysis of $m_k$ separates the two ways a sign can arise — real/rational pole structure (compatible with non-reversible finite models) versus a complex pair at irrational phase \emph{on the dominant circle}, which by the theory of nonnegatively realizable sequences excludes \emph{every} finite-state model; the dominance qualifier is load-bearing (subdominant complex pairs are unconstrained) and the protocol tests it explicitly.  The pole stage requires \emph{order selection} (numerical Hankel rank before Prony) — fixed-order fitting provably hallucinates poles, as we show on an exactly solvable model.  Applied to Google's 2023 surface-code memory experiment (public data; 64 stabilizer streams, 50{,}000 shots each): \textbf{64/64 streams nominally sign} — decisively for the typical stream ($\lambda_{\min} \sim -10^{-3}$, tight bootstrap intervals at $B = 2000$, surviving a conservative Bonferroni statement, an order of magnitude above the derived within-shot drift systematic), \textbf{provisionally for the weakest streams} (the 9 streams with $|\lambda_{\min}| < 5\times 10^{-4}$, where the residual concern is the bootstrap's boundary behavior — not drift, which Prop. 2.2 shows is second-order and PSD-directed, i.e. it can only mask signs, never mimic them); the order-selected pole stage finds \textbf{purely real primary poles in 64/64 streams}, and the rank scan finds \textbf{no stream with a resample-stable, dominant complex pair} — the clock signature, which the oscillatory control exhibits at stability and dominance $1.00$.  The verdict is anchored by an exact model lemma: the standard detection-event mechanism — measurement errors firing the same detector in consecutive rounds — is a four-state chain of letter-transfer rank two with exact poles $\{+0.2836, -0.2336\}$; it signs the offset Hankel with purely real poles, while being \emph{distributionally time-reversal invariant} and \emph{Mealy-realizable on a reversible chain} — so the sign certifies the absence of a reversible \textbf{Moore-form} hidden realization, not any temporal arrow in the data and not the absence of all reversible-chain models, and we phrase all conclusions accordingly.  Conclusion: the device's syndrome statistics show exactly the pair-mechanism structure its source publication documents by other means, and no rotation ("clock") signature that would obstruct finite-state modeling up to the dimension bound finite data can certify. Controls demonstrate instrument power: an i.i.d. stream passes (the degenerate-Hankel bootstrap bias is corrected by basic-bootstrap intervals), and a hidden irrational rotation at matched rate is decisively rejected.

\section*{1. Introduction}
Decoders and noise analyses for quantum error correction presuppose noise models — usually finite-state, often implicitly reversible in their statistical structure.  Parametric checks (fitting a candidate model and inspecting residuals) cannot exclude the \emph{class}: a poor fit indicts the candidate, not the family.  Structural tests can: they certify properties that every member of a model class must imprint on the data, so a violation excludes the class wholesale, at all dimensions.

The structural property we use is a moment condition on run statistics.  Its correct statement turned out to be subtler than the natural first guess, and Section 2 develops both the theorem and the counterexample that disciplines it — a correction we highlight because the naive version is tempting, simple, and wrong in a way that real data punishes (Section 4's device results \emph{reverse} under the correction).

\section*{2. Theory}
\subsection*{2.1 The sign stage}
Let $s_t \in \{0,1\}$ be stationary, $p(1^n) = \Pr[s_t = \cdots = s_{t+n-1} = 1]$, and $m_k := p(1^{k+1})$ for $k \ge 0$.

\textbf{Proposition 2.1 (Moore form).}  If $s$ is generated by a finite hidden-Markov model \textbf{in Moore form} — emissions conditioned on the state, $\Pr[s_t = 1 \mid X_t = i] = e_i$ — whose chain $T$ is reversible with respect to its stationary law $\pi$ (detailed balance), then $(m_k)_{k\ge 0}$ is a Hamburger moment sequence; in particular every offset Hankel $H_m = [m_{i+j}]_{i,j=0}^m$ and every doubly-offset $[m_{i+j+2}]$ is PSD.

\emph{Proof.}  Let $E = \operatorname{diag}(e_i)$ and $D = \operatorname{diag}(\pi)$.  Then $m_k = \mathbf 1^{\mathsf T} E D\,(TE)^k\,\mathbf 1$. Reversibility makes $K = D^{1/2} T D^{-1/2}$ symmetric; with $u = E^{1/2} D^{1/2}\mathbf 1$ and $M = E^{1/2} K E^{1/2}$ (symmetric), a direct computation gives $m_k = \langle u, M^k u\rangle$ — first for invertible $E$, then for general $E \succeq 0$ by continuity ($E \mapsto E + \epsilon I$, $\epsilon \downarrow 0$; the bilinear quantity $\langle u_\epsilon, M_\epsilon^k u_\epsilon\rangle$ is polynomial in $\epsilon$, as is $m_k$ of the perturbed model).  By the spectral theorem $m_k = \int x^k\, d\mu(x)$ with $\mu$ the positive spectral measure of $M$ at $u$.  $\square$

\textbf{The Moore qualifier cannot be dropped (scope correction).}  An earlier draft claimed the Mealy (edge-emission) case follows "by converting to Moore form on the edge set, with reversibility inherited" — that is \textbf{false}, and we record the refutation because the error is natural: the induced Moore model lives on the \emph{edge chain}, whose back-transition $(j,k) \to (i,j)$ requires $k = i$, so detailed balance fails identically for any non-degenerate chain.  The counterexample is this paper's own Lemma 2.3: the detection-event process is a Mealy model on a \emph{reversible} chain ($\mu_t$ i.i.d. — detailed balance trivially — with symmetric edge emissions $\Pr[s = 1 \mid \mu_{t-1}, \mu_t]$ depending on the transition), it reproduces the signed process exactly, and random search shows the phenomenon is generic ($1449/2000$ random Mealy models on reversible chains sign the offset Hankel, worst $\lambda_{\min} = -1.26\times 10^{-2}$ — numerically close to, but distinct from, Lemma 2.3's exact $-1.32\times 10^{-2}$; fixed-seed receipt).  The certified class of the sign stage is therefore exactly: \emph{Moore-form HMMs with detailed-balance chains}.

\textbf{The trap (why the offset matters).}  The naive Hankel $[p(1^{i+j})]_{i,j\ge 0}$ \emph{includes} $p(1^0) = 1$, which is not part of the moment structure.  Counterexample: the deterministic alternator (two states, $A \to B \to A$; $A$ emits $1$, $B$ emits $0$) is a reversible two-state model with $p(1) = 1/2$, $p(11) = 0$; the naive Hankel $\begin{psmallmatrix}1 & 1/2\\ 1/2 & 0\end{psmallmatrix}$ is indefinite — the naive test rejects a legitimate finite model.  The offset Hankel is exactly $\begin{psmallmatrix}1/2 & 0\\ 0 & 0\end{psmallmatrix} \succeq 0$, as Proposition 2.1 demands (here $M$ is nilpotent and $\mu = \tfrac12\delta_0$).

\textbf{Proposition 2.2 (across-shot robustness — and its limits).}  A convex mixture of Hamburger sequences is Hamburger.  Hence pooling over \emph{shots} — ensemble drift, parameter wander between repetitions — cannot produce a false sign.  \textbf{This does not cover within-shot temporal drift}: the windowed run estimator pools over time, and the pooled run statistics of a nonstationary stream are not a convex mixture of per-model moment sequences.  \textbf{The mechanism, corrected in review} (an earlier draft derived a first-order effect $+\mu^3 b$ from the $k$-shortened anchor window; that derivation forgot that the run's \emph{interior} samples advance the rate by $+bk/2$ on average, exactly cancelling the window's $-bk/2$): for an i.i.d. stream with linearly drifting rate $p_t = a + bt$, the windowed estimator's mean rate is \emph{independent of the moment order at first order}, so

\[
m_1^2 - m_0 m_2 \;=\; O(b^2),
\]
and the leading second-order term is the mixture-variance $-\mu^2\,\mathrm{Var}(p_t) < 0$ — \textbf{toward PSD, for either drift sign} (verified exactly: $-4.3\times 10^{-6}$ at $\mu = 0.15$, $b = +0.002$/round, $T = 25$; same sign for $b < 0$).  Linear within-shot drift is therefore a false-\textbf{pass} systematic, not a false-sign one: it cannot mimic a SIGNED verdict at leading order, which \emph{strengthens} the device results below.  Magnitude: with $\Delta$ the measured within-shot range of the round means ($0.002$–$0.028$ across streams — a range, not a per-round slope), $\mu^2\Delta^2/12 \lesssim 1.5\times 10^{-6}$: three orders below typical signs.  The residual weak-stream caveat in Section 4 is therefore carried by the bootstrap boundary issue alone, not by drift; per-round-block detrended analysis remains the named upgrade for non-linear drift shapes.

\subsection*{2.2 The pole stage}
A signed offset Hankel excludes reversible finite models.  Two mechanisms produce signs:

\textbf{(a) Real/rational structure.}  Non-reversible finite models can violate the Cauchy–Schwarz relation $m_1^2 \le m_0 m_2$ — strong pair correlation with suppressed triples — while keeping all poles of $m_k$ real (including negative) or rational-phase.  These remain finite-state models.

\textbf{(b) Irrational rotation.}  If $m_k$ carries a complex pole pair at irrational phase \textbf{on its dominant circle}, no finite-state model reproduces it: the precise statement is the Berstel–Soittola theory of nonnegatively realizable ($\mathbb R_+$-rational) series — the dominant singularities of such a sequence are the radius times roots of unity [Berstel–Reutenauer Ch. 8, which treats the $\mathbb R_+$ case our application needs, not only $\mathbb N$] — which also covers the residue subtlety that bare Perron–Frobenius does not (peripheral eigenvalues of the matrix can cancel from $m_k$; the theorem speaks of the sequence's own singularities).  The theorem constrains \emph{dominant} singularities only: a \textbf{subdominant} complex pair at any phase is fully compatible with finite nonnegative realization, so the protocol must — and now does — test dominance explicitly: the discovery criterion requires the resampled ratio $|z_{\mathrm{cx}}|/\hat\rho$ (pair modulus over fitted spectral radius) to be consistent with $1$ within its own tolerance, folded into the $N(\delta)$ quantifier below.  Two honest qualifications: (i) a phase estimated to precision $\pm\delta$ is consistent with a rational $p/q$ for $q \gtrsim \pi/\delta$, so finite data excludes finite-state models only \emph{up to a dimension} $N(\delta)$ — every "excludes all models" below is to be read with that quantifier; (ii) pooled data with parameter drift is not exactly finite-pole (a continuous mixture of geometric decays has a cut), so the pole fit is an effective low-order description and must carry order selection and uncertainty, as follows.

Stage two, order-selected: estimate the numerical rank $r$ of the moment Hankel from its singular-value profile, then solve the order-$r$ Prony system.  \textbf{This is not optional}: fitting fixed order $m = 3$ to a rank-2 sequence hallucinates a pole.  On the exactly-solvable model below, with \emph{exact} moments and the least-squares Prony estimator of our implementation, forced order-3 reports a spurious third pole at $-0.046$ where the exact spectrum is $\{+0.2836, -0.2336\}$; on \emph{estimated} moments the artifact's value wanders with the noise and the estimator variant (an earlier draft quoted $0.175$ from a simulation run — the very fact that the number is not stable across estimators is the point), while rank selection recovers the exact pair in all variants. Real or rational-phase poles: mechanism (a).  A stable irrational pair on the dominant circle, robust to rank choice and resampling: mechanism (b).

\textbf{Lemma 2.3 (the standard mechanism signs, with real poles — exact).}  The standard detection-event model of the QEC literature [Dennis et al. 2002; the source experiment's own correlation analysis] — $s_t = \mu_t \oplus \mu_{t-1} \oplus d_t$ with i.i.d. measurement errors $\mu$ (rate $q$) and data errors $d$ (rate $r$): measurement errors fire the same detector in consecutive rounds — violates $m_1^2 \le m_0 m_2$ and signs the offset Hankel, with purely real poles.  This is \emph{exactly computable}: the process is a four-state chain in $(\mu_{t-1}, \mu_t)$ whose letter-transfer has rank two; at $(q, r) = (0.08, 0.05)$ the exact moments give $m_1^2 - m_0 m_2 = +2.786\times 10^{-3}$, $\lambda_{\min} = -1.3\times 10^{-2}$, and exact poles $\{+0.28358,\ -0.23358\}$ — the negative real pole is the fingerprint of the pair structure.  (The chain is \emph{not} reversible as a hidden realization even though the process is invariant in distribution under time reversal — see the caution below.)

\textbf{Caution on the word "non-reversible" — two distinctions, not one.}  The model above is distributionally time-reversal invariant ($p(\tilde w) = p(w)$: relabel $\nu_t = \mu_{-t}$), yet it signs; and it \emph{is} realizable as a Mealy model on a reversible chain (Section 2.1), yet it signs. What the sign stage detects is therefore the absence of a \emph{reversible Moore-form hidden realization} — a representation-class property — \textbf{not} temporal irreversibility of the data, and \textbf{not} the absence of every reversible-chain model; direct word-reversal statistics (e.g. $p(011)$ vs $p(110)$) are expected to pass and do. We phrase all verdicts accordingly.  A second consequence, stated plainly: since the detection-event construction (XOR of consecutive measurements) guarantees the pair mechanism whenever measurement errors exist, the sign stage \emph{should} fire on essentially every QEC detector stream; its role here is validation that the instrument reads known structure correctly from raw run statistics.  The discovery channel is the pole stage.

\section*{3. Statistical procedure}
\begin{enumerate}
\item Estimate $p(1^n)$, $n = 1, \ldots, 2m{+}1$, by the windowed run estimator over bulk rounds and all shots ($m = 3$ here).
\item Form the offset Hankel; compute $\lambda_{\min}$.
\item Bootstrap over shots ($B = 2000$; per-shot run-count resampling, identical statistic).  Exactly singular Hankels (e.g. i.i.d.: $m_k = p^{k+1}$, rank one) bias the plug-in $\lambda_{\min}$ negative; the \textbf{basic bootstrap} interval $[2\hat\lambda - q_{0.975},\, 2\hat\lambda - q_{0.025}]$ corrects the \emph{first-order bias}.  Stated precisely: at the rank-deficient boundary the bootstrap is \textbf{not consistent} for extremal eigenvalues [Andrews 2000] — the basic interval does not restore consistency, it passes the controls; for weak streams ($|\lambda_{\min}| < 5\times 10^{-4}$; 9 of 64) verdicts are labeled provisional, and $m$-out-of-$n$ subsampling (which has known boundary validity) is the named procedure for them. Verdict SIGNED only if the corrected upper limit is below zero.
\item Pole stage: order-selected Prony on $(m_k)$; report imaginary fractions, phases, and — whenever a complex pair appears — the dominance ratio $|z_{\mathrm{cx}}|/\hat\rho$ with its resampled interval.
\item \textbf{Rank scan} (added after the conservative primary selection was found to \emph{under}-select on the oscillatory control, whose moment Hankel has $\sigma_2/\sigma_1 \approx 6\times 10^{-3}$, below the $2\times 10^{-2}$ tolerance): at every order $r \le m$, fit poles on the full sample; whenever a complex pair appears, refit at that order on each bootstrap moment vector and report \emph{stability} (fraction of resamples showing a pair), median dominance, and median phase.  A clock candidate requires a pair that is \textbf{both} resample-stable (stability $\ge 0.9$) \emph{and} dominant (median $|z_{\mathrm{cx}}|/\hat\rho \ge 0.95$) — thresholds fixed here, before the data column is read; the two failure modes this rules out being order-misselection artifacts (unstable) and subdominant pairs (unconstrained by the realizability theorem).  For calibration: the oscillatory control scores stability $1.00$, dominance $1.00$; the two device streams with any stable pair score $(0.68, 0.83)$ and $(0.90, 0.55)$ — each failing one criterion.
\end{enumerate}
\section*{4. Results on public surface-code data}
Google 2023 surface-code memory experiment (Zenodo 6804040): detection events parsed from the stim \texttt{.b8} files, detector coordinates from circuit files; five distance-3 patch configurations (four patch centers, one of them in both bases) and distance-5, 25 rounds, 50{,}000 shots per configuration; bulk rounds.

\begin{small}\begin{verbatim}
64 stabilizer streams (5 distance-3 patch configurations x 8
stabilizers = 40, plus distance-5 x 24); B = 2000:
  event rates 0.10-0.21;  round-to-round drift 0.002-0.028
  SIGN STAGE: 64/64 SIGNED (no reversible Moore-form HMM)
    typical lam_min ~ -1e-3, tight CIs (stream (5.0,6.0) of the
    bZ_d3_r25_center_3_5 patch, lam_min -1.99e-3, shown as the
    example: [-2.1e-3, -1.9e-3]);
    Bonferroni (64 tests, 5% family, x1.6 margin inflation,
    normal approx): ALL SIGNED verdicts survive (worst inflated
    upper bound -1.7e-5);
    the drift systematic of Prop. 2.2 is second-order, PSD-
    directed, and <~ 1.5e-6: it cannot mimic any sign; the
    provisional label below concerns the bootstrap boundary
    only (9 streams have |lam_min| < 5e-4)
  POLE STAGE (order-selected): 64/64 purely real primary poles
    (51 streams at numerical rank 1, 13 at rank 2; max
    |Im z|/|z| = 0.000).  Honesty note: a literal rank-1
    nonnegative sequence is automatically Hamburger, so the
    MECHANISM confirmation (negative pole at rank 2) rests on
    the 13 rank-2 streams; the 51 rank-1 readings reflect the
    conservative 2e-2 rank tolerance, the same under-selection
    mode the control exposes (the scan stage covers it).  The fixed-order version of this stage
    had reported 2/64 complex fits; both were order-misselection
    artifacts and vanish under rank selection - documented
    because it is exactly the failure mode Lemma 2.3's exact
    model exposes.
  RANK SCAN: 2/64 streams show a resample-stable complex pair at
    some order <= 3 - and both fail dominance, exactly as the
    theory sorts them: one at stability 0.68, dominance 0.83,
    with only TWO length-7 runs feeding its top moment; one at
    stability 0.90 but dominance 0.55 (clearly subdominant =
    unconstrained).  0/64 streams show a stable AND dominant
    pair - the clock signature.
  RUN-COUNT DIAGNOSTIC: length-7 runs per stream min 2, median
    96, max 399 - the min-2 stream is the stability-0.68 scan
    entry above, and its top moments deserve no confidence.
controls (matched rate and sample size):
  iid:        lam_min -1.5e-05, CI [-2.9e-05, +5.3e-05] -> PASS
              (rank scan: none stable)
  irrational
  rotation:   lam_min -9.4e-04, CI [-1.0e-03, -8.5e-04] -> SIGNED
              rank scan: stable pair at r=2 (stability 0.90,
              dominance 1.00) and r=3 (stability 1.00, dominance
              1.00) - the stable-AND-dominant clock signature
              the device data does not show.  NOTE the honest
              instrument subtlety this control exposed: the
              conservative PRIMARY order selection (tolerance
              2e-2) under-selects here (reads rank 1, real pole)
              because the pair carries ~0.6% of the Hankel
              weight; the scan stage exists precisely for this,
              and an earlier draft's claim that the primary
              order-selected analysis alone finds the control's
              pair was wrong and is corrected by the scan.
\end{verbatim}\end{small}
\textbf{Reading, with the right verbs.}  The device's syndrome statistics admit \emph{no reversible Moore-form hidden realization} — with the same real-pole structure (negative pole present at rank 2) as the standard pair mechanism of Lemma 2.3, whose presence in this dataset is independently documented by the source experiment's own correlation analysis (which we thereby confirm from run statistics alone, a different and cruder observable).  The pole stage finds \textbf{no stable dominant complex signature}: no clock, and nothing obstructing finite-state noise modeling of these devices up to the dimension bound of Section 2.2(i).  Magnitude caveat, stated: Lemma 2.3's illustrative $(q, r)$ gives $\lambda_{\min}$ an order larger than the typical device value; we claim mechanism agreement (sign + real-pole structure + negative pole), not a per-stream quantitative fit, which would require matching $(q, r)$ per stabilizer — a worthwhile follow-up.

\textbf{The cautionary tale, quantified.}  Under the naive (unoffset) statistic these same 64 streams all \emph{pass} — the weaker, wrongly justified test sees nothing.  The correction simultaneously fixed the false-rejection trap (the alternator) and exposed the real structure that the naive version was blind to; the pole stage then needed its own correction (order selection) to stop inventing structure.  We document both because all the versions are natural, only one chain is theorem-backed, and they disagree on real data.

\textbf{Statistical caveats, collected.}  (i) The i.i.d. null sits at a boundary (rank-one Hankel) where bootstrap inference for extremal eigenvalues is known to be delicate [Andrews 2000]; the basic bootstrap correction ($B = 2000$) is motivated and passes the controls but does \emph{not} restore boundary consistency, and a full size/power study over the composite null (reversible models near rank deficiency, matched rates), plus $m$-out-of-$n$ subsampling for the weak streams, is the named upgrade.  (ii) The 64 streams share shots and hardware and are spatially correlated; the conservative Bonferroni statement above ($\times 1.6$ margin inflation, all 64 surviving) controls the family-wise error under arbitrary dependence (it is a union bound — dependence only makes it conservative); the genuine residual fragility is the per-test bootstrap normal approximation near the PSD boundary, flagged under (i).  (iii) Top moments at $m = 3$ rest on few long runs per stream — now published per stream (min $2$, median $96$, max $399$); the min-$2$ stream's scan entry is dismissed on exactly this ground, and the rank selection partially insulates the pole stage elsewhere.

\section*{5. What a discovery would look like}
A SIGNED stream whose rank scan shows a \emph{stable} complex pair — reproducible across patch centers, robust to windowing and to rank choice, with a resampled phase interval excluding the low-order rationals — \textbf{and whose resampled dominance ratio $|z_{\mathrm{cx}}|/\hat\rho$ is consistent with $1$} (the condition without which the realizability theorem says nothing) would certify that no finite-state noise model below the corresponding dimension bound $N(\delta)$ reproduces that stream, with consequences for decoder noise priors.  The present data shows the opposite, which is itself useful: leakage, TLS defects, and drift on these devices are of the kind hidden finite-state structure can express [McEwen et al.; Chen et al.].

Natural extensions: composite-block run statistics (patterns rather than the symbol $1$) sharpen the moment conditions; cross-stabilizer joint statistics; per-segment analysis around documented burst events.

\section*{Reproducibility}
A single script parses the public data and regenerates every number above (fixed seeds; bit-identical reruns), including the controls and the textbook-model simulation.

\section*{References}
\begin{itemize}
\item Google Quantum AI, Suppressing quantum errors by scaling a surface code logical qubit, \emph{Nature} \textbf{614} (2023) 676; data: Zenodo 6804040 — including the supplementary detection-event correlation ($p_{ij}$) analysis our Section 4 confirms from run statistics.
\item E. Dennis, A. Kitaev, A. Landahl, J. Preskill, Topological quantum memory, \emph{J. Math. Phys.} \textbf{43} (2002) 4452 — the standard detection-event/pair model.
\item J. Berstel, C. Reutenauer, \emph{Noncommutative Rational Series with Applications}, CUP (2011), Ch. 8 — dominant singularities of positively realizable series; the chapter treats the $\mathbb R_+$-rational case our application needs (via Soittola's theorem), not only $\mathbb N$-rational (Stage 2's theorem).
\item A. Heller, \emph{Ann. Math. Statist.} \textbf{36} (1965) 1286 — cone characterization of finite functions of Markov chains (the classical frame for Section 2's realization claims).
\item D. W. K. Andrews, Inconsistency of the bootstrap when a parameter is on the boundary, \emph{Econometrica} \textbf{68} (2000) 399 — the boundary caveat of Section 4.
\item B. Efron, R. Tibshirani, \emph{An Introduction to the Bootstrap} — basic bootstrap intervals.
\item M. McEwen et al., Resolving catastrophic error bursts from cosmic rays, \emph{Nat. Phys.} (2022); Z. Chen et al., Exponential suppression of bit or phase errors, \emph{Nature} \textbf{595} (2021) — device temporal-correlation context.
\item Reversible-chain spectral background: J. F. C. Kingman/D. G. Kendall tradition; C. J. Geyer, Practical MCMC, \emph{Statist. Sci.} \textbf{7} (1992) 473 — initial-sequence estimators exploit exactly Prop. 2.1's structure.
\end{itemize}
\end{document}
